\chapter{研究成果与分析（核心章节）}

\section{主要研究结果}
项目获得了具有代表性的高温红外屏蔽复合样品。与对照组相比，优化样品在目标波段表现出更低的平均发射率，并在高温氧化后仍保持较高结构稳定性。实验表明，适度提高致密相比例并构建多尺度界面，可有效抑制高温辐射增强效应。

相关文献对比可见表~\ref{tab:EAB}。

\begin{table}[H]
\centering
\caption{关键性能对比表:}
\begin{tabularx}{0.95\textwidth}{C{2.2cm}C{3.2cm}C{3.2cm}C{3.2cm}}
\toprule
组别 & 8--14\,$\mu$m 平均发射率（900\,$^{\circ}$C） & 900\,$^{\circ}$C 氧化后质量保持率（\%） & 抗热震循环次数（次） \\
\midrule
实验组 & 0.41 & 96.8 & 25 \\
对照组 & 0.57 & 91.2 & 14 \\
\bottomrule
\end{tabularx}
\end{table}

\section{学术论文成果（重点）}
\noindent\textbf{论文信息：}
\begin{itemize}
  \item 标题：面向极端热环境的高温红外屏蔽复合材料设计与性能优化
  \item 期刊名称：Ceramics International
  \item 发表状态：已发表
  \item 发表时间：2025年12月15日
  \item 期刊等级：中科院2区，JcrQ1
  \item 作者排序：项目负责人（第一作者），指导教师（通讯作者）
  \item 项目资助标注：本研究由广东省自然科学基金(2023A1515030106)、国家自然科 学基金(52272058)、广州市科技计划项目基础与应用基础研究专题 (2023A04J1006)、广州市教育局材料科学与工程重点学科( 202255464)、广州大学本科生创新创业训练计划(XJ202511078412)以及广州大学 研究生创新训练计划(JCCX2025054)提供资助。
\end{itemize}

\noindent\textbf{论文核心内容提炼：}
论文围绕“高温稳定性与红外屏蔽性能协同提升”这一核心问题，构建了复合材料组分设计与工艺优化策略。通过对不同配比样品在中高温条件下的红外发射率、氧化增重与微观相演变进行系统表征，揭示了界面致密化程度与辐射特性之间的关联机制。研究结果表明，优化后的多相复合结构可在 900\,$^{\circ}$C 附近维持较低波段发射率，同时显著减缓氧化导致的性能衰退。论文进一步提出了面向工程应用的参数选取区间，为高温热管理和红外抑制材料开发提供了可复用的实验路径与数据依据。
